\section{Light --- Lib --- Merge Sort}
%%% generated by the blueprint pipeline from TCSlib.Programs.Light.Lib.MergeSort
\subsection*{Light --- Lib --- Merge Sort: Introduction}

Declarations of the Lean module \texttt{TCSlib.Programs.Light.Lib.MergeSort}.

\subsection*{Light --- Lib --- Merge Sort: Results}

\begin{definition}[Merge-sort length constant]
\lean{Light.MergeSort.Len}
\label{Light.MergeSort.Len}
\leanok
The merge-sort length constant is $0$.
\end{definition}

\begin{definition}[Merge-sort start constant]
\lean{Light.MergeSort.Start}
\label{Light.MergeSort.Start}
\leanok
\uses{Light.MergeSort.Len}
The merge-sort start constant is $1$.
\end{definition}

\begin{definition}[Merge-sort free-space constant]
\lean{Light.MergeSort.Free}
\label{Light.MergeSort.Free}
\leanok
\uses{Light.MergeSort.Len}
The merge-sort free-space constant is $2$.
\end{definition}

\begin{definition}[Merge-sort half constant]
\lean{Light.MergeSort.Half}
\label{Light.MergeSort.Half}
\leanok
\uses{Light.MergeSort.Len}
The merge-sort half constant is $3$.
\end{definition}

\begin{definition}[Length and unused variable]
\lean{Light.MergeSort.Unused}
\label{Light.MergeSort.Unused}
\leanok
\uses{Light.MergeSort.Len}
The abbreviation for the length is $0$, and the abbreviation for the unused variable
index is $4$.
\end{definition}

\begin{definition}[Merge and copy step]
\lean{Light.sortJoin}
\label{Light.sortJoin}
\leanok
\uses{Light.MergeSort.Free, Light.MergeSort.Half, Light.MergeSort.Len, Light.MergeSort.Start, Light.MergeSort.Unused, Light.Stmt, Light.v}
For procedure indices $p_{\mathrm{merge}}$ and $p_{\mathrm{copy}}$, define a statement
that calls the merge procedure with arguments $\mathrm{Start}$, $\mathrm{Half}$,
$\mathrm{Start}+\mathrm{Half}$, the truncated difference of $\mathrm{Len}$ and
$\mathrm{Half}$, and $\mathrm{Free}$, assigning the result to $\mathrm{Unused}$. It then
calls the copy procedure with arguments $\mathrm{Free}$, $\mathrm{Start}$, and
$\mathrm{Len}$, assigning that result to $\mathrm{Unused}$ as well.
\end{definition}

\begin{definition}[Recursive sorting body]
\lean{Light.sortBody}
\label{Light.sortBody}
\leanok
\uses{Light.MergeSort.Free, Light.MergeSort.Half, Light.MergeSort.Len, Light.MergeSort.Start, Light.MergeSort.Unused, Light.Stmt, Light.k, Light.sortJoin, Light.v}
For natural-number procedure indices
$p_{\mathrm{sort}},p_{\mathrm{half}},p_{\mathrm{merge}},p_{\mathrm{copy}}$, define a
statement that tests whether $1<\mathrm{Len}$. When this holds, it calls the halving
procedure with argument $\mathrm{Len}$, calls the sorting procedure on the two segments
with arguments $(\mathrm{Half},\mathrm{Start},\mathrm{Free})$ and
$(\mathrm{Len}-\mathrm{Half},\mathrm{Start}+\mathrm{Half},\mathrm{Free})$, respectively,
and then performs the merge and copy statement. When $1<\mathrm{Len}$ does not hold, it
does nothing.
\end{definition}

\begin{definition}[Sorting time bound]
\lean{Light.sortTime}
\label{Light.sortTime}
\leanok
For each natural number $n$, the specified time is $61n\lceil\log_2 n\rceil+74(n-1)+4$,
where $\lceil\log_2 n\rceil$ denotes the natural-number ceiling logarithm returned by
$\operatorname{clog}(2,n)$.
\end{definition}

\begin{theorem}[Ceiling logarithm under a split]
\lean{Light.MergeSort.clog_split}
\label{Light.MergeSort.clog_split}
\leanok
\difficulty{4}
For every natural number $n\ge2$, the following bound holds, where
$\operatorname{clog}_2(m)$ is the least natural exponent $k$ for which $m\le2^k$:
\[\left\lfloor\frac{n+1}{2}\right\rfloor\operatorname{clog}_2\!\left(\left\lfloor\frac{n+1}{2}\right\rfloor\right)+\left(n-\left\lfloor\frac{n+1}{2}\right\rfloor\right)\operatorname{clog}_2\!\left(n-\left\lfloor\frac{n+1}{2}\right\rfloor\right)+n\le n\operatorname{clog}_2(n).\]
\end{theorem}

\begin{definition}[Merge sort procedure indices]
\lean{Light.MergeSort.Procs}
\label{Light.MergeSort.Procs}
\leanok
\uses{Light.Program, Light.copyBody, Light.halfBody, Light.mergeBody, Light.sortBody}
A program $P$ together with four natural-number indices
$p_{\mathrm{sort}},p_{\mathrm{half}},p_{\mathrm{merge}},p_{\mathrm{copy}}$ has the
required procedure layout when its entries at those indices are, respectively, the
sorting body parameterized by these four indices, the halving body, the merging body,
and the copying body.
\end{definition}

\begin{theorem}[Merge sort time split bound]
\lean{Light.MergeSort.time_split}
\label{Light.MergeSort.time_split}
\leanok
\uses{Light.MergeSort.clog_split, Light.sortTime}
\difficulty{2}
Define $T(n)=61n\lceil\log_2 n\rceil+74(n-1)+4$, where the logarithmic factor is the
least natural number $k$ such that $n\le 2^k$. For every natural number $n\ge2$,
$$61n+70+T\!\left(\left\lfloor\frac{n+1}{2}\right\rfloor\right)+T\!\left(n-\left\lfloor\frac{n+1}{2}\right\rfloor\right)\le
T(n).$$
\end{theorem}

\begin{definition}[Sorted segment with preserved cells]
\lean{Light.MergeSort.Sorted}
\label{Light.MergeSort.Sorted}
\leanok
\uses{Light.KeptBut, Light.Seg}
For functions $\mu,\mu':\mathbb{N}\to\mathbb{Z}$, natural numbers $a,fr$, and a list $L$
of integers, the property holds when both of the following are true: there is a list
$L'$ occupying the segment specified by $\mu'$ starting at $a$, such that $L'$ is a
permutation of $L$ and is nondecreasing; and the memory specified by $\mu'$ is obtained
from $\mu$ by preserving the required cells outside that segment for a length of $L$.
\end{definition}

\begin{theorem}[Reflexive sortedness for short lists]
\lean{Light.MergeSort.sorted_refl}
\label{Light.MergeSort.sorted_refl}
\leanok
\uses{Light.MergeSort.Sorted, Light.SameOn.refl, Light.Seg}
\difficulty{4}
Let $\mu:\mathbb{N}\to\mathbb{Z}$ be a memory function, and let $L$ be an integer list
of length at most one. If $\mu$ stores $L$ starting at address $a$, then $L$ is sorted
at $a$ below $fr$ relative to $\mu$: a permutation of $L$ is stored there in
nondecreasing order, and $\mu$ agrees with itself below $fr$ outside the interval
starting at $a$ of length $|L|$.
\end{theorem}

\begin{theorem}[Sorted merge preserves sortedness]
\lean{Light.MergeSort.sorted_merge}
\label{Light.MergeSort.sorted_merge}
\leanok
\uses{Light.KeptBut, Light.MergeSort.Sorted, Light.SameOutside, Light.Seg, List.Pairwise.merge}
\difficulty{5}
Let $\mu,\mu_1,\mu_2,\mu_3,\mu_4$ be integer-valued memory functions, and let
$L,S_1,S_2$ be integer lists. Choose natural numbers $a,fr,h,m$ with $h\le |L|$. Suppose
$S_1$ is a permutation of the first $h$ entries of $L$ and is nondecreasing, and that
changing $\mu$ to $\mu_1$ preserves all cells below $fr$ outside the interval starting
at $a$ whose length is the length of that prefix. Suppose $S_2$ is a permutation of the
remaining entries of $L$ and is nondecreasing, and changing $\mu_1$ to $\mu_2$ preserves
all cells below $fr$ outside the interval starting at $a+h$ whose length is the length
of that suffix. Suppose $\mu_2$ and $\mu_3$ agree outside the interval starting at $fr$
of length $m$, the memory $\mu_4$ stores the merge of $S_1$ and $S_2$ starting at $a$,
and $\mu_3$ and $\mu_4$ agree outside the interval starting at $a$ of length $|L|$. Then
$\mu_4$ witnesses that $L$ is sorted at $a$ below $fr$ relative to the original memory
$\mu$: a permutation of $L$ is stored there in nondecreasing order, and the required
cells below $fr$ outside that interval are preserved.
\end{theorem}

\begin{theorem}[Merge sort join postcondition]
\lean{Light.MergeSort.join_ends}
\label{Light.MergeSort.join_ends}
\leanok
\uses{Light.Ends, Light.MergeSort.Procs, Light.SameOutside, Light.Seg, Light.copy_meets, Light.frame, Light.merge_meets, Light.sortJoin}
\difficulty{4}
Let $P$ contain the designated sorting, halving, merging, and copying procedures, and
suppose the word bound is at least the space bound. Let $d,a,fr,h,n$ be natural numbers,
let $r$ be an integer, and let $S_1,S_2$ be integer lists of lengths $h$ and $n-h$,
respectively, with $h\le n$ and $1\le n$. Suppose the memory function $\mu$ stores $S_1$
starting at $a$ and $S_2$ starting at $a+h$, with $a+n\le fr$ and
$fr+n\le\mathrm{space}$, and suppose $d<\mathrm{depth}$. Then the join procedure,
started with frame $(n,a,fr,h,r)$ and memory $\mu$, has an execution of cost at most
$56n+34$ ending in a state with memory $\sigma'.\mathrm{mem}$ for which there is a
memory function $\mu'$ that agrees with $\mu$ outside the interval starting at $fr$ of
length $n$, stores the merge of $S_1$ and $S_2$ starting at $a$, and agrees with
$\sigma'.\mathrm{mem}$ outside the interval starting at $a$ of length $n$.
\end{theorem}

\begin{theorem}[Merge sort body runtime]
\lean{Light.MergeSort.body_ends}
\label{Light.MergeSort.body_ends}
\leanok
\uses{Light.Ends, Light.Ends.iteLast, Light.Ends.mono, Light.Ends.skip, Light.Meets, Light.MergeSort.Procs, Light.MergeSort.Sorted, Light.MergeSort.join_ends, Light.MergeSort.sorted_merge, Light.MergeSort.sorted_refl, Light.MergeSort.time_split, Light.Seg, Light.Seg.congr, Light.Seg.drop, Light.Seg.take, Light.frame, Light.half_meets, Light.sortBody, Light.sortTime}
\difficulty{6}
Let a program contain the prescribed sorting, halving, merging, and copying procedures,
and let the space limit be at most the word limit, with one also at most the word limit.
Let $L$ be a list of integers represented in memory on the segment beginning at $a$,
with $a+|L|\le fr$ and $fr+|L|$ at most the space limit. Suppose $d$ is below the depth
limit, and that for every starting address $a'$, list $L'$, and memory $\mu'$, if
$1<|L|$, $2|L'|\le |L|+1$, $L'$ is represented in $\mu'$ on the segment beginning at
$a'$, and $a'+|L'|\le fr$, then the sorting procedure at depth $d+1$, started with
length $|L'|$, address $a'$, and free address $fr$, terminates within
$\mathrm{sortTime}(|L'|)$ and leaves memory in which the segment beginning at $a'$ is
sorted and the required memory outside that segment is preserved. Then the sorting body
at depth $d$, started with length $|L|$, address $a$, free address $fr$, and memory
$\mu$, terminates within $\mathrm{sortTime}(|L|)$ and leaves memory in which the segment
beginning at $a$ is sorted and the required memory outside it is preserved.
\end{theorem}

\begin{theorem}[Merge sort correctness and resource bound]
\lean{Light.sort_meets}
\label{Light.sort_meets}
\leanok
\uses{Light.Meets, Light.Meets.of_body, Light.MergeSort.Procs, Light.MergeSort.Sorted, Light.MergeSort.body_ends, Light.Seg, Light.sortTime}
\difficulty{6}
Let $lim$ be resource limits, and let $P$ contain the sort, halving, merge, and copy procedures.
Suppose $lim.space\le lim.word$ and $1\le lim.word$. Let $L$ be a list of integers,
$a$ its starting address, $fr$ a scratch-space address, and $\mu$ a memory storing $L$
at $a$. Assume $a+|L|\le fr$ and $fr+|L|\le lim.space$. For every recursion depth $d$
with $d+\mathrm{clog}_2(|L|)+1\le lim.depth$, sorting the $|L|$ items at $a$ using
$fr$ terminates within
$61|L|\,\mathrm{clog}_2(|L|)+74(|L|-1)+4$ steps (with natural-number subtraction),
and leaves at $a$ a sorted permutation of $L$. Among addresses below $fr$, only the
interval from $a$ of length $|L|$ may have changed.
\end{theorem}
